\chapter[Lezione 14 --- 27 ottobre 2026]{Lezione 14: operazioni sui limiti e limiti unilateri}
\label{chap:lezione-14}

\section{Teoremi sulle operazioni dei limiti}
\label{sec:lez14-operazioni-limiti}

Si considerano i limiti delle operazioni tra funzioni
introdotte nella \cref{sec:lez07-operazioni}.

\subsection{Limite della somma}
\label{sec:lez14-limite-somma}

\begin{theorem}[Limite della somma]
\label{thm:lez14-limite-somma}
Siano \(f,g\colon A\longrightarrow\R\), con
\(A\subseteq\R\) non vuoto, e sia
\(x_0\in\overline{\R}\) un punto di accumulazione per \(A\).
Si supponga che
\[
  \lim_{x\to x_0}f(x)=\ell_1\in\overline{\R},
  \qquad
  \lim_{x\to x_0}g(x)=\ell_2\in\overline{\R}.
\]
Allora:
\begin{enumerate}
  \item se \(\ell_1,\ell_2\in\R\), si ha
  \[
    \lim_{x\to x_0}(f+g)(x)=\ell_1+\ell_2;
  \]
  \item se almeno uno dei due limiti è \(+\infty\) e
  nessuno è \(-\infty\), si ha
  \[
    \lim_{x\to x_0}(f+g)(x)=+\infty;
  \]
  \item se almeno uno dei due limiti è \(-\infty\) e
  nessuno è \(+\infty\), si ha
  \[
    \lim_{x\to x_0}(f+g)(x)=-\infty.
  \]
\end{enumerate}
\end{theorem}

\begin{proof}
\emph{Primo punto.}
Sia \(\varepsilon>0\). Per la \cref{def:lez08-limite},
esistono intorni \(U'_{x_0}\) e \(U''_{x_0}\) tali che
\[
  |f(x)-\ell_1|<\frac{\varepsilon}{2}
  \qquad\forall x\in(U'_{x_0}\cap A)\setminus\{x_0\},
\]
\[
  |g(x)-\ell_2|<\frac{\varepsilon}{2}
  \qquad\forall x\in(U''_{x_0}\cap A)\setminus\{x_0\}.
\]
Per la \cref{sec:lez11-intersezione-intorni}, l'insieme
\[
  U_{x_0}:=U'_{x_0}\cap U''_{x_0}
\]
è un intorno di \(x_0\). Per ogni
\(x\in(U_{x_0}\cap A)\setminus\{x_0\}\), dalla disuguaglianza
triangolare segue
\[
\begin{aligned}
  |(f+g)(x)-(\ell_1+\ell_2)|
    &=|(f(x)-\ell_1)+(g(x)-\ell_2)|\\
    &\le |f(x)-\ell_1|+|g(x)-\ell_2|\\
    &<\frac{\varepsilon}{2}+\frac{\varepsilon}{2}
     =\varepsilon.
\end{aligned}
\]
Per la \cref{def:lez08-limite}, segue la tesi.

\emph{Secondo punto.}
Per la commutatività della somma, si può supporre
\(\ell_1=+\infty\).
Se \(\ell_2\in\R\), fissato \(K\in\R\), esistono
intorni \(U'_{x_0}\) e \(U''_{x_0}\) nei quali valgono,
rispettivamente,
\[
  f(x)>K-\ell_2+1,
  \qquad g(x)>\ell_2-1.
\]
Nell'intersezione dei due intorni si ha
\[
  (f+g)(x)>K.
\]
Se invece \(\ell_2=+\infty\), si scelgono intorni nei
quali \(f(x)>K\) e \(g(x)>0\); sulla loro intersezione
vale nuovamente \((f+g)(x)>K\).
In entrambi i casi le disuguaglianze valgono nei punti
di \(A\) distinti da \(x_0\). Poiché \(K\) è arbitrario,
segue il limite \(+\infty\).

\emph{Terzo punto.}
Per commutatività, si supponga \(\ell_1=-\infty\).
Se \(\ell_2\in\R\), fissato \(K\in\R\), si scelgono
intorni nei quali
\[
  f(x)<K-\ell_2-1,
  \qquad g(x)<\ell_2+1.
\]
Sulla loro intersezione si ottiene \((f+g)(x)<K\).
Se \(\ell_2=-\infty\), si scelgono intorni nei quali
\(f(x)<K\) e \(g(x)<0\), ottenendo la stessa disuguaglianza.
La \cref{def:lez08-limite} fornisce il limite \(-\infty\).
\end{proof}

\begin{example}[Forma indeterminata \(+\infty-\infty\)]
\label{ex:lez14-somma-indeterminata}
Per \(x\to+\infty\), si considerino tre coppie di funzioni:
\[
\begin{aligned}
  f(x)=x,\quad g(x)=-x
    &\quad\Longrightarrow\quad (f+g)(x)=0,\\
  f(x)=2x,\quad g(x)=-x
    &\quad\Longrightarrow\quad (f+g)(x)=x,\\
  f(x)=\frac{x}{2},\quad g(x)=-x
    &\quad\Longrightarrow\quad (f+g)(x)=-\frac{x}{2}.
\end{aligned}
\]
In tutti e tre i casi \(f(x)\to+\infty\) e
\(g(x)\to-\infty\), mentre la somma tende rispettivamente
a \(0\), \(+\infty\) e \(-\infty\).
\end{example}

\begin{corollary}[Somma con una funzione localmente limitata da un lato]
\label{cor:lez14-somma-localmente-limitata}
Siano \(f,g\colon A\longrightarrow\R\), con
\(A\subseteq\R\) non vuoto, e sia
\(x_0\in\overline{\R}\) un punto di accumulazione per \(A\).
Allora:
\begin{enumerate}
  \item se \(\displaystyle\lim_{x\to x_0}f(x)=+\infty\)
  ed esiste un intorno \(U_{x_0}\) tale che
  \[
    \inf\{g(x):x\in(U_{x_0}\cap A)\setminus\{x_0\}\}
    \in\R,
  \]
  allora
  \[
    \lim_{x\to x_0}(f+g)(x)=+\infty;
  \]

  \item se \(\displaystyle\lim_{x\to x_0}f(x)=-\infty\)
  ed esiste un intorno \(U_{x_0}\) tale che
  \[
    \sup\{g(x):x\in(U_{x_0}\cap A)\setminus\{x_0\}\}
    \in\R,
  \]
  allora
  \[
    \lim_{x\to x_0}(f+g)(x)=-\infty.
  \]
\end{enumerate}
\end{corollary}

La dimostrazione del \cref{cor:lez14-somma-localmente-limitata}
non viene sviluppata in questa sede.

\begin{example}[Somma con una funzione limitata]
\label{ex:lez14-somma-seno}
La funzione \(\sin x\) è limitata su \(\R\), poiché
\(-1\le\sin x\le1\) per ogni \(x\in\R\).
La limitatezza sull'intero dominio implica sia la
limitatezza inferiore sia quella superiore in ogni intorno.
Per il \cref{cor:lez14-somma-localmente-limitata}, si ottiene
\[
  \lim_{x\to+\infty}(x+\sin x)=+\infty,
  \qquad
  \lim_{x\to-\infty}(x+\sin x)=-\infty.
\]
Il primo limite è stato verificato anche mediante confronto
nell'\cref{ex:lez12-limite-x-seno}.
\end{example}

\subsection{Limite del prodotto}
\label{sec:lez14-limite-prodotto}

\begin{theorem}[Limite del prodotto]
\label{thm:lez14-limite-prodotto}
Siano \(f,g\colon A\longrightarrow\R\), con
\(A\subseteq\R\) non vuoto, e sia
\(x_0\in\overline{\R}\) un punto di accumulazione per \(A\).
Si supponga che
\[
  \lim_{x\to x_0}f(x)=\ell_1\in\overline{\R},
  \qquad
  \lim_{x\to x_0}g(x)=\ell_2\in\overline{\R}.
\]
Allora:
\begin{enumerate}
  \item se \(\ell_1,\ell_2\in\R\), si ha
  \[
    \lim_{x\to x_0}(fg)(x)=\ell_1\ell_2;
  \]
  \item se \(\ell_1\in(0,+\infty)\cup\{+\infty\}\) e
  \(\ell_2=+\infty\), rispettivamente \(\ell_2=-\infty\),
  si ha
  \[
    \lim_{x\to x_0}(fg)(x)=+\infty,
    \qquad\text{rispettivamente }-\infty;
  \]
  \item se \(\ell_1\in(-\infty,0)\cup\{-\infty\}\) e
  \(\ell_2=+\infty\), rispettivamente \(\ell_2=-\infty\),
  si ha
  \[
    \lim_{x\to x_0}(fg)(x)=-\infty,
    \qquad\text{rispettivamente }+\infty.
  \]
\end{enumerate}
Le affermazioni valgono anche scambiando \(f\) e \(g\).
\end{theorem}

La dimostrazione del \cref{thm:lez14-limite-prodotto}
non viene sviluppata in questa sede.

\begin{example}[Forma indeterminata \(0\cdot\infty\)]
\label{ex:lez14-prodotto-indeterminato}
Per \(x\to+\infty\), si considerino le funzioni
\[
\begin{aligned}
  f(x)=\frac1x,\quad g(x)=x
    &\quad\Longrightarrow\quad (fg)(x)=1,\\
  f(x)=\frac1{x^2},\quad g(x)=x
    &\quad\Longrightarrow\quad (fg)(x)=\frac1x,\\
  f(x)=\frac1x,\quad g(x)=x^2
    &\quad\Longrightarrow\quad (fg)(x)=x.
\end{aligned}
\]
In tutti e tre i casi \(f(x)\to0\) e \(g(x)\to+\infty\),
mentre il prodotto tende rispettivamente a \(1\), \(0\)
e \(+\infty\).
\end{example}

\subsection{Limite del rapporto}
\label{sec:lez14-limite-rapporto}

\begin{theorem}[Limite del rapporto]
\label{thm:lez14-limite-rapporto}
Siano \(f,g\colon A\longrightarrow\R\), con
\(A\subseteq\R\) non vuoto, e sia
\(x_0\in\overline{\R}\) un punto di accumulazione per \(A\).
Si supponga che \(g(x)\ne0\) nei punti di \(A\) distinti
da \(x_0\) di un opportuno intorno di \(x_0\) e che
\[
  \lim_{x\to x_0}f(x)=\ell_1\in\overline{\R},
  \qquad
  \lim_{x\to x_0}g(x)=\ell_2\in\overline{\R}.
\]
Allora:
\begin{enumerate}
  \item se \(\ell_1\in\R\) e \(\ell_2\in\R\setminus\{0\}\),
  si ha
  \[
    \lim_{x\to x_0}\frac{f(x)}{g(x)}=\frac{\ell_1}{\ell_2};
  \]
  \item se \(\ell_1\in\R\) e
  \(\ell_2\in\{-\infty,+\infty\}\), si ha
  \[
    \lim_{x\to x_0}\frac{f(x)}{g(x)}=0;
  \]
  \item se \(\ell_1=+\infty\) e \(\ell_2\in\R\setminus\{0\}\),
  il rapporto tende a \(+\infty\) per \(\ell_2>0\) e a
  \(-\infty\) per \(\ell_2<0\); se \(\ell_1=-\infty\),
  i due risultati si scambiano;
  \item se \(\ell_1\in\{+\infty,-\infty\}\) e \(\ell_2=0\),
  e \(g\) ha segno costante in un intorno bucato di \(x_0\),
  il rapporto tende a \(+\infty\) quando i segni di
  \(f\) e \(g\) coincidono e a \(-\infty\) quando sono
  opposti;
  \item se \(\ell_1\in\R\setminus\{0\}\) e \(\ell_2=0\),
  e \(g\) ha segno costante in un intorno bucato di \(x_0\),
  vale la stessa conclusione del punto precedente.
\end{enumerate}
\end{theorem}

La dimostrazione del \cref{thm:lez14-limite-rapporto}
non viene sviluppata in questa sede.

\begin{example}[Forme indeterminate del rapporto]
\label{ex:lez14-rapporto-indeterminato}
Per \(x\to+\infty\), le coppie
\[
\begin{aligned}
  f(x)=\frac1x,\quad g(x)=\frac1x
    &\quad\Longrightarrow\quad \frac{f(x)}{g(x)}=1,\\
  f(x)=\frac1{x^2},\quad g(x)=\frac1x
    &\quad\Longrightarrow\quad \frac{f(x)}{g(x)}=\frac1x,\\
  f(x)=\frac1x,\quad g(x)=\frac1{x^2}
    &\quad\Longrightarrow\quad \frac{f(x)}{g(x)}=x
\end{aligned}
\]
forniscono la forma \(0/0\), ma rapporti con limiti
rispettivamente \(1\), \(0\) e \(+\infty\).
Analogamente, le coppie
\[
\begin{aligned}
  f(x)=x,\quad g(x)=x
    &\quad\Longrightarrow\quad \frac{f(x)}{g(x)}=1,\\
  f(x)=x,\quad g(x)=x^2
    &\quad\Longrightarrow\quad \frac{f(x)}{g(x)}=\frac1x,\\
  f(x)=x^2,\quad g(x)=x
    &\quad\Longrightarrow\quad \frac{f(x)}{g(x)}=x
\end{aligned}
\]
mostrano che anche la forma \(+\infty/+\infty\) è
indeterminata.
\end{example}

\section{Avvicinamento al limite dall'alto e dal basso}
\label{sec:lez14-avvicinamento-valore}

\begin{definition}[Avvicinamento dall'alto e dal basso]
\label{def:lez14-avvicinamento-valore}
Siano \(f\colon A\longrightarrow\R\), con
\(A\subseteq\R\) non vuoto, e
\(x_0\in\overline{\R}\) un punto di accumulazione per \(A\).
Si supponga che
\[
  \lim_{x\to x_0}f(x)=\ell\in\R.
\]
Si dice che \(f(x)\) \emph{tende a \(\ell\) dall'alto} se
esiste un intorno \(U_{x_0}\) tale che
\[
  f(x)>\ell
  \qquad\forall x\in(U_{x_0}\cap A)\setminus\{x_0\}.
\]
Analogamente, si dice che \(f(x)\) \emph{tende a
\(\ell\) dal basso} se, nello stesso senso,
\(f(x)<\ell\). Questi comportamenti si indicano
rispettivamente con \(f(x)\to\ell^+\) e
\(f(x)\to\ell^-\) per \(x\to x_0\).
\end{definition}

Le scritture \(\ell^+\) e \(\ell^-\) specificano il verso
di avvicinamento al valore \(\ell\): non rappresentano
ulteriori numeri ai quali applicare le regole algebriche
dei limiti.

\clearpage

\begin{example}[Verso di avvicinamento di una somma]
\label{ex:lez14-verso-somma}
Siano
\[
  f(x):=3-x^2,
  \qquad g(x):=2+x^4.
\]
Per \(x\to0\), si ha \(f(x)\to3\) dal basso e
\(g(x)\to2\) dall'alto. Tuttavia,
\[
  (f+g)(x)=5+x^4-x^2<5
  \qquad\text{se }0<|x|<1.
\]
Il limite della somma è \(5\), raggiunto dal basso:
il verso di avvicinamento non si determina sommando
formalmente i simboli \(3^-\) e \(2^+\).
\end{example}

\section{Limiti destri e sinistri}
\label{sec:lez14-limiti-unilateri}

Il verso di avvicinamento al valore limite della
\cref{sec:lez14-avvicinamento-valore} è distinto dal lato
da cui la variabile \(x\) si avvicina al punto \(x_0\).

\begin{definition}[Intorni destri e sinistri]
\label{def:lez14-intorni-unilateri}
Siano \(x_0\in\R\) e \(\delta>0\). A partire dall'intorno
circolare della \cref{def:lez07-intorno}, si definiscono
l'\emph{intorno destro} e l'\emph{intorno sinistro}
rispettivamente mediante
\[
\begin{aligned}
  D^+(x_0,\delta)
    &:=D(x_0,\delta)\cap\{x\in\R:x>x_0\}
      =(x_0,x_0+\delta),\\
  D^-(x_0,\delta)
    &:=D(x_0,\delta)\cap\{x\in\R:x<x_0\}
      =(x_0-\delta,x_0).
\end{aligned}
\]
\end{definition}

\begin{definition}[Punti di accumulazione destri e sinistri]
\label{def:lez14-accumulazione-unilaterale}
Sia \(A\subseteq\R\) non vuoto e sia \(x_0\in\R\).
Il punto \(x_0\) si dice \emph{di accumulazione destro}
per \(A\) se
\[
  D^+(x_0,\delta)\cap A\ne\varnothing
  \qquad\forall\delta>0.
\]
Si dice \emph{di accumulazione sinistro} per \(A\) se
\[
  D^-(x_0,\delta)\cap A\ne\varnothing
  \qquad\forall\delta>0.
\]
\end{definition}

\begin{example}[Accumulazione agli estremi di un intervallo]
\label{ex:lez14-accumulazione-unilaterale}
Per \(A=(3,4)\), il punto \(3\) è di accumulazione
destro, ma non sinistro, e \(4\) è di accumulazione
sinistro, ma non destro. Ogni \(x_0\in(3,4)\) è di
accumulazione sia destro sia sinistro.
\end{example}

\begin{figure}[htbp]
  \centering
  \begin{tikzpicture}[unipd/diagram,x=2.25cm,y=1cm]
  \draw[->,unipd-muted] (2.5,0) -- (4.52,0) node[right] {\(\R\)};
  \draw[very thick,unipd-primary] (3,0) -- (4,0);
  \filldraw[fill=white,draw=unipd-primary,line width=.85pt] (3,0) circle (2.15pt);
  \filldraw[fill=white,draw=unipd-primary,line width=.85pt] (4,0) circle (2.15pt);
  \fill[unipd-primary] (3.5,0) circle (1.75pt);
  \node[below=5pt] at (3,0) {\(3\)};
  \node[below=5pt] at (3.5,0) {\(x_0\)};
  \node[below=5pt] at (4,0) {\(4\)};
  \node[above=14pt] at (3.5,0) {\(A=(3,4)\)};
  \draw[->,unipd-accent,line width=.75pt] (3.05,.35) -- (3.29,.35);
  \draw[->,unipd-accent,line width=.75pt] (3.95,.35) -- (3.71,.35);
\end{tikzpicture}

  \caption{Per \(A=(3,4)\), il punto \(3\) è di
  accumulazione destro, \(4\) è di accumulazione
  sinistro e ogni \(x_0\in A\) è di entrambi i tipi.}
  \label{fig:lez14-accumulazione-unilaterale}
\end{figure}

\clearpage

\begin{definition}[Limite destro e limite sinistro]
\label{def:lez14-limiti-unilateri}
Siano \(f\colon A\longrightarrow\R\), con
\(A\subseteq\R\) non vuoto, e \(\ell\in\overline{\R}\).
Se \(x_0\in\R\) è di accumulazione destro per \(A\),
si dice che \(f\) \emph{ammette limite destro} \(\ell\)
per \(x\to x_0\), e si scrive
\[
  \lim_{x\to x_0^+}f(x)=\ell,
\]
se, per ogni intorno \(V_\ell\) di \(\ell\), esiste
\(\delta>0\) tale che
\[
  f(x)\in V_\ell
  \qquad\forall x\in D^+(x_0,\delta)\cap A.
\]
Se \(x_0\) è di accumulazione sinistro, il
\emph{limite sinistro} è definito analogamente,
sostituendo \(D^+\) con \(D^-\) e scrivendo
\(\lim_{x\to x_0^-}f(x)=\ell\).
\end{definition}

\begin{theorem}[Limite e limiti unilateri]
\label{thm:lez14-limite-unilateri}
Siano \(f\colon A\longrightarrow\R\), con
\(A\subseteq\R\) non vuoto, e \(x_0\in\R\) un punto
di accumulazione sia destro sia sinistro per \(A\).
Per ogni \(\ell\in\overline{\R}\), vale
\[
  \lim_{x\to x_0}f(x)=\ell
  \quad\Longleftrightarrow\quad
  \lim_{x\to x_0^-}f(x)
  =\lim_{x\to x_0^+}f(x)=\ell.
\]
\end{theorem}

La dimostrazione del \cref{thm:lez14-limite-unilateri}
non viene sviluppata in questa sede.
